% ============================================================
% Come si usano, argomento 14: variabili continue multiple
% ============================================================
\argomento{Variabili aleatorie continue multiple}{Formule a p.~\pageref*{f:congiuntac}}[Variabili continue
  multiple]

\begin{scelta}{Come scegliere la formula}
\renewcommand{\arraystretch}{1.25}
\begin{tabular}{@{}P{0.55\linewidth}P{0.39\linewidth}@{}}
  \toprule
  \textbf{Se il testo\ldots} & \textbf{si usa\ldots} \\
  \midrule
  chiede la probabilità che la coppia cada in una regione & l'integrale doppio della densità \vv{congiuntac},
    a fette \\
  chiede la densità di una sola variabile & le marginali \vv{marginalic} \\
  chiede se le variabili sono indipendenti & la fattorizzazione \vv{indipc} \\
  chiede la densità di $X$ sapendo $Y = y$, o di risalire da $Y$ a $X$ & \vv{condc}, Bayes \vv{totalec} \\
  chiede la media di $g(X, Y)$ & \vv{mediac2}, o in due tempi \vv{mediacondc} \\
  chiede covarianza o correlazione & \vv{covc} \\
  dà la CDF congiunta, o chiede la probabilità di un rettangolo & \vv{cdfcongiunta} \\
  \bottomrule
\end{tabular}

Tutto si traduce dal discreto (argomenti 8 e 9) con gli integrali al posto delle somme. Prima di integrare conviene
disegnare la regione: gli estremi dell'integrale interno possono dipendere dalla variabile esterna.
\end{scelta}

\scheda{congiuntac}{
  \usoquando{Si chiede la probabilità di un evento che coinvolge entrambe le variabili: $X + Y \leq 1$, $X < Y$, il
    punto in un cerchio.}
  \usocome{Si disegna la regione nel piano, si sceglie quale variabile tenere fissa e si integra a fette
    (\vv{integraledoppio}).}
  \usoesempio{$X$ e $Y$ esponenziali indipendenti di parametro 1: $P(X + Y \leq 1) = \int_0^1 \int_0^{1-x} e^{-x}
    e^{-y} \dd y \dd x = 1 - \frac{2}{e} \simeq 0{,}264$.}
}

\scheda{marginalic}{
  \usoquando{Si conosce la densità congiunta e serve quella di una sola variabile.}
  \usocome{Si integra via la variabile che non interessa, su tutto il suo supporto.}
  \usoesempio{$f_{X,Y}(x, y) = x + y$ su $[0, 1]^2$: $f_X(x) = \int_0^1 (x + y) \dd y = x + \frac{1}{2}$, e per
    simmetria $f_Y(y) = y + \frac{1}{2}$.}
}

\scheda{indipc}{
  \usoquando{Per verificare se due variabili sono indipendenti, o per costruire la congiunta di variabili
    indipendenti (prodotto delle densità).}
  \usoattenzione{Basta un punto in cui la congiunta non è il prodotto delle marginali.}
  \usoesempio{$x + y$ su $[0, 1]^2$: in $(0, 0)$ la congiunta vale 0 mentre $f_X(0)\, f_Y(0) = \frac{1}{4}$: non sono
    indipendenti. Due esponenziali indipendenti: $f_{X,Y} = e^{-x} e^{-y}$ per $x, y \geq 0$.}
}

\scheda{condc}{
  \usoquando{Si vuole la densità di $X$ sapendo che $Y$ ha assunto il valore $y$.}
  \usocome{Si prende la «fetta» della congiunta a $y$ fissato e la si divide per la sua area, che è $f_Y(y)$.}
  \usoattenzione{$P(Y = y) = 0$: la condizionata si ottiene condizionando a «$Y$ vicino a $y$», ed è il rapporto fra
    densità.}
  \usoesempio{$x + y$ su $[0, 1]^2$: $f_{X \mid Y}(x \mid y) = \frac{x + y}{y + 1/2}$; con $y = 0$ vale $2x$, con
    $y = 1$ vale $\frac{x + 1}{1{,}5}$.}
}

\scheda{totalec}{
  \usoquando{Si conosce una densità marginale e il «meccanismo» $f_{Y \mid X}$, e serve l'altra marginale o la
    condizionata al contrario: sono la probabilità totale e Bayes con le densità.}
  \usocome{Come nel discreto, con l'integrale al posto della somma.}
  \usoesempio{$x + y$ su $[0, 1]^2$: $\int_0^1 f_{X \mid Y}(x \mid y)\, f_Y(y) \dd y = \int_0^1 (x + y) \dd y =
    x + \frac{1}{2} = f_X(x)$.}
}

\scheda{mediac2}{
  \usoquando{Si vuole la media di una combinazione delle due variabili: $XY$, $X^2$, $X + Y$.}
  \usocome{Integrale doppio di $g$ pesata con la congiunta. Se $g$ è lineare basta la linearità.}
  \usoesempio{$x + y$ su $[0, 1]^2$: $E[XY] = \int_0^1 \int_0^1 (x^2 y + x y^2) \dd y \dd x = \frac{1}{3}
    \cdot \frac{1}{2} + \frac{1}{2} \cdot \frac{1}{3} = \frac{1}{3}$.}
}

\scheda{mediacondc}{
  \usoquando{La media è facile sapendo il valore di una delle due variabili.}
  \usocome{Prima la media condizionata, come funzione di $y$; poi la sua media rispetto a $f_Y$.}
  \usoesempio{$x + y$ su $[0, 1]^2$: $E[X \mid Y = y] = \frac{1/3 + y/2}{y + 1/2}$ ($\frac{2}{3}$ per $y = 0$,
    $\frac{5}{9}$ per $y = 1$); poi $E[X] = \int_0^1 \left(\frac{1}{3} + \frac{y}{2}\right) \dd y = \frac{7}{12}$.}
}

\scheda{covc}{
  \usoquando{Si vuole sapere quanto due variabili continue si muovono insieme.}
  \usocome{$E[XY]$ con l'integrale doppio, le medie e le varianze dalle marginali, poi le formule del discreto.}
  \usoesempio{$x + y$ su $[0, 1]^2$: $COV = \frac{1}{3} - \left(\frac{7}{12}\right)^2 = -\frac{1}{144}$,
    $\sigma_X^2 = \sigma_Y^2 = \frac{11}{144}$, quindi $\rho = -\frac{1}{11} \simeq -0{,}09$: correlazione negativa ma
    debole.}
}

\scheda{cdfcongiunta}{
  \usoquando{Probabilità del tipo «$X \leq x$ e $Y \leq y$», probabilità di rettangoli, o il passaggio dalla CDF
    congiunta alla densità.}
  \usocome{Per un rettangolo si usano i quattro vertici: più quello in alto a destra e quello in basso a sinistra,
    meno gli altri due. Dalla CDF alla densità si deriva una volta per ogni variabile (\vv{derivateparziali}).}
  \usoesempio{$x + y$ su $[0, 1]^2$: $F_{X,Y}(x, y) = \frac{x^2 y + x y^2}{2}$ sul quadrato;
    $P\left(\frac{1}{2} < X \leq 1, \frac{1}{2} < Y \leq 1\right) = 1 - \frac{3}{8} - \frac{3}{8} + \frac{1}{8} =
    \frac{3}{8}$.}
}
